%%%%%%%%%%%%%%%%%%%%%%%%%%%%%%%%%%%%%%%%%%%%%%%%%%%%%%%%%%%%%%%%%%%%%%
% Lớp 10 | Thống kê | Số gần đúng. Sai số
% Dạng 0D6?1-3: Quy tròn số gần đúng
% Số câu: 21  (sinh tự động bởi tach_cau.py; từ nay thầy sửa trực tiếp file này)
%%%%%%%%%%%%%%%%%%%%%%%%%%%%%%%%%%%%%%%%%%%%%%%%%%%%%%%%%%%%%%%%%%%%%%

%%%%%%%%%%%%%%%% TRẮC NGHIỆM 4 PHƯƠNG ÁN %%%%%%%%%%%%%%%%

% nguồn: MaHoaNH_TracNghiem_Lop10
\begin{ex}%[0D6N1-3]%[ma:76cd7182]
	Số quy tròn đến hàng  phần nghìn của số $ a=0{,}1234$ là
	 \choice
		{$ 0{,}124$} 
		{$ 0{,}12$} 
		{\True $  0{,}123$}
		{$ 0{,}13$} 
	 \loigiai{
	$  a=0{,}1234 \approx 0{,}123$.}
\end{ex}

% nguồn: MaHoaNH_TracNghiem_Lop10
\begin{ex}%[0D6N1-3]%[ma:2e61c0ba]
	Số quy tròn của $20 \, 182 \, 020$  đến hàng trăm là
	\choice 
		 {\True $20\, 182 \, 000$}
		{$20 \, 180  \, 000$}
		{$20\, 182  \, 100$}
		{$20\,  182\, 020$}
	 \loigiai{$20 \, 182 \, 020$ $\approx 20\, 182 \, 000$.
	 }
\end{ex}

% nguồn: MaHoaNH_TracNghiem_Lop10
\begin{ex}%[0D6N1-3]%[ma:da64f228]
	Khi sử dụng máy tính bỏ túi với $ 10$ chữ số thập phân ta được $\sqrt{8}=2{,}828427125$. Giá trị gần đúng của $\sqrt{8}$ chính xác đến hàng phần trăm  là
	 \choice
		{$ 2{,}81$} 
		{\True  $ 2{,}83$}
		{$  2{,}82$}
		{$  2{,}80$}
	 \loigiai{
		 $2,828427125 \approx 2{,}83$.
	 }
\end{ex}

% nguồn: MaHoaNH_TracNghiem_Lop10
\begin{ex}%[0D6N1-3]%[ma:f8506026]
	Khi sử dụng máy tính bỏ túi với $10$ chữ số thập phân ta được $\sqrt{8}=2{,}828427125$. Giá trị gần đúng của $\sqrt{8}$  chính xác đến hàng phần trăm là
	 \choice
		{$2{,}81$} 
		{\True  $2{,}83$}
		 {$2{,}82$}
		 {$2{,}80$}
	 \loigiai{
		Giá trị gần đúng của $\sqrt{8}=2{,}828427125$ chính xác đến hàng phần trăm  là $2{,}83$.
	 }
\end{ex}

% nguồn: MaHoaNH_TracNghiem_Lop10
\begin{ex}%[0D6N1-3]%[ma:d7eb2e32]
	Số quy tròn của số $5{,}1472$ với  độ chính xác $d=0{,}05$ là
	 \choice
		{$5$} 
		{$5{,}15$} 
		{\True  $5{,}1$}
		 {$5{,}2$}
	\loigiai{ 
	Vì độ chính xác $d=0{,}05$ đến hàng phần trăm, nên ta  làm tròn số $a=5{,}1472$ đến hàng phần chục. Vậy số quy tròn của $a$ là $5{,}1$.
	} 
\end{ex}

% nguồn: MaHoaNH_TracNghiem_Lop10
\begin{ex}%[0D6H1-3]%[ma:9289100d]
	Cho số $a=367\,653\,964\pm 213$.  Số quy tròn của số gần đúng $367\,653\,964$ là
	\choice 
		{$367\,653\,960$} 
		{$36\,7653\,000$} 
		{\True  $367\,654\,000$}
		{$367\,653\,970$} 
	 \loigiai{
	Vì độ chính xác đến hàng trăm $d=213$ nên số quy tròn của số gần đúng $367\,653\,964$  là $36\,7654\,000$.}
\end{ex}

% nguồn: MaHoaNH_TracNghiem_Lop10
\begin{ex}%[0D6H1-3]%[ma:ec2016c8]
	Theo thống kê, dân số Việt Nam năm $2016$ được ghi lại như sau $\overline{S}=94\,444\,200\pm 3000$ (người). Số quy tròn của số gần đúng  $94\,444\,200$ là
	\choice 
		{\True  $94\,440\,000$}
		{$94\,450\,000$} 
		{$94\,444\,000$} 
		{$94\,400\,000$} 
	\loigiai{ 
	Vì $1000<3000<10000$ nên hàng cao nhất mà $d$ nhỏ hơn một đơn vị của hàng đó là hàng chục nghìn. Nên ta phải quy tròn số $94\,444\,200$ đến  hàng chục nghìn. Vậy số quy tròn là $94\,440\,000$.}
\end{ex}

% nguồn: MaHoaNH_TracNghiem_Lop10
\begin{ex}%[0D6H1-3]%[ma:67890935]
	Cho  $\overline{a}=31\,462\,689\pm 150$. Số quy tròn của số $31\,462\,689$ là
	\choice 
		{$31\,462\,000$} 
		{$31\,463\,700$} 
		{$31\,463\,600$} 
		 {\True $31\,463\,000$}
	\loigiai{ 
	Độ chính xác đến hàng trăm $(d=150)$ nên ta quy tròn đến hàng nghìn. Vậy số quy tròn của số  $31\,462\,689$ là $31\,463\,000$.}
\end{ex}

% nguồn: MaHoaNH_TracNghiem_Lop10
\begin{ex}%[0D6H1-3]%[ma:e72c929a]
	Cho giá trị gần đúng của $\pi$ là $a=3,141\,592\,653\,589$ với độ chính xác $10^{-10}$ ($10$ chữ số thập phân). Số quy tròn của  $a$ là
	\choice 
		 {\True $a=3{,}141592654$}
		 {$a=3{,}1415926536$}
		 {$a=3{,}141592653$}
		 {$a=3{,}1415926535$}
	 \loigiai{
		Ta có $10^{-11}<10^{-10}<10^{-9}$ nên hàng cao nhất mà $d$ nhỏ hơn một đơn vị của hàng đó là hàng phần tỉ. Do đó ta phải quy tròn số $a=3{,}141\,592\,653\,589$  đến hàng phần tỉ.\\
	Vậy,  số quy tròn là $a= 3{,}141\,592\,654$.}
\end{ex}

% nguồn: MaHoaNH_TracNghiem_Lop10
\begin{ex}%[0D6H1-3]%[ma:a628ba95]
	Cho số $\overline{a}=17\,658\pm 16$.  Số quy tròn của số gần đúng $17\,658$ là
	 \choice
		 {$18\,000$}
		{$17\,800$} 
		 {$17\,600$}
		 {\True $17\,700$}
	\loigiai{ 
		Ta có $10<16<100$ nên hàng cao  nhất mà $d$ nhỏ hơn một đơn vị của hàng đó là hàng trăm. Do đó ta phải quy tròn số $176\,38$ đến hàng trăm.\\
		Vậy số quy tròn là  $17\,700$ (hay viết $\overline{a}\approx 17\,700$).
	 }
\end{ex}

% nguồn: MaHoaNH_TracNghiem_Lop10
\begin{ex}%[0D6H1-3]%[ma:bdf5fe08]
	Số tiền quỹ lớp $10$A còn lại là $a=1\,647\,500$ (đồng) với độ chính xác $d=500$ (đồng). Hãy viết  số quy tròn của số $a$?
	\choice 
		{\True  Số quy tròn của $a$ là $1\,648\,000$ (đồng)}
		{Số quy tròn của $a$ là $1\,647\,000$ (đồng)} 
		{Số quy tròn của $a$ là $1\,649\,000$  (đồng)}
		{Số quy tròn của $a$ là $1\,650\,000$  (đồng)}
	\loigiai{ 
		Vì độ chính xác $d=500$ đến hàng trăm nên ta quy tròn số gần đúng $a$ đến  hàng nghìn.\\
		Vậy, số quy tròn của $a$  là $1\,648\,000$ (đồng).
	} 
\end{ex}

% nguồn: MaHoaNH_TracNghiem_Lop10
\begin{ex}%[0D6H1-3]%[ma:a45b29a0]
	Chiều dài gần đúng của một cái bàn học là $a=1{,}238$ m với độ chính xác $d=0{,}01$ m. Hãy  viết số quy tròn của số $a$?
	 \choice
		 {\True Số quy tròn của $a$ là $1{,}2$ m}
		{Số quy tròn của  $a$ là $1{,}23$ m}
		{Số quy tròn của $a$  là $1{,}24$ m}
		{Số  quy tròn của $a$ là $1{,}248$ m}
	\loigiai{ 
		Vì độ chính xác $d=0{,}01$ đến hàng phần trăm nên ta quy tròn số gần đúng $a$ đến hàng  phần chục.\\
		Vậy, số quy tròn của $a$  là $1{,}2$ m.
	} 
\end{ex}

% nguồn: MaHoaNH_TracNghiem_Lop10
\begin{ex}%[0D6H1-3]%[ma:19276cd4]
	Theo thống kê, dân số Việt Nam năm $2002$ là $79\,715\,675$ người. Giả sử sai số tuyệt đối của số  liệu thống kê này nhỏ hơn $10\,000$ người. Hãy viết số quy tròn của số trên
	 \choice
		 {$79\,710\,000$ người}
		{$79\,716\,000$ người} 
		{\True $79\,720\,000$  người}
		 {$79\,700\,000$ người}
	\loigiai{ 
		Vì sai số tuyệt đối của  số liệu thống kê này nhỏ hơn $10\,000$ người nên độ chính xác đến hàng nghìn nên ta quy tròn đến hàng chục nghìn.\\
		Vậy số quy tròn cùa số trên là  $79\,720\,000$ người.
	 }
\end{ex}

% nguồn: MaHoaNH_TracNghiem_Lop10
\begin{ex}%[0D6H1-3]%[ma:c1dc03e4]
	Độ dài các cạnh của một đám vườn hình chữ nhật là $x=7{,}8$ $\pm 2$ cm và $y=25{,}6$ m $\pm 4$ cm. Cách viết chuẩn của diện tích (sau khi quy  tròn) là
	 \choice
		{$200\mathrm{~m}^2\pm 0{,}9\mathrm{~m}^2$} 
		 {$199\mathrm{~m}^2\pm 0{,}8\mathrm{~m}^2$}
		 {$199\mathrm{~m}^2\pm 1\mathrm{~m}^2$}
		{\True $200\mathrm{~m}^2\pm 1\mathrm{~m}^2$} 
	\loigiai{ 
		 Ta có
		\begin{eqnarray*}
			x=7{,}8 \text{ m} \pm 2 \text{ cm} &\Rightarrow& 7{,}78 \text{ m}\leq x\leq  7{,}82 \text{ m}.\\
			y=25{,}6 \text{ m} \pm  4 \text{ cm} &\Rightarrow& 25{,}56 \text{ m} \leq y\leq 25{,}64 \text{ m}.
		\end{eqnarray*}
		Do  đó diện tích của hình chữ nhật thỏa mãn $198{,}8568\mathrm{~m}^2\leq xy\leq 200{,}5048\mathrm{~m}^2$.\\
		Vậy, cách viết chuẩn của diện tích sau khi quy tròn là $200$ m$^2\pm 1$  m$^2$.
	} 
\end{ex}

% nguồn: Ok_TracNghiemLop11_Lan1
\begin{ex}%[0D6K1-3]%[ma:b8c4f0b1]
	Trên đường tròn lượng giác gốc $A$, cung lượng giác nào có các điểm biểu  diễn tạo thành tam giác đều?
	\choice 
		{$\dfrac{k\pi}{3}$} 
		{$k\pi  $}
		 {\True $\dfrac{k2\pi}{3}$}
		 {$\dfrac{k\pi}{2}$}
	 \loigiai{
		Tam giác đều có góc ở đỉnh là $60^{\circ}$ nên góc ở tâm là $120^{\circ}$ tương ứng  $\dfrac{k2\pi}{3}$.
	 }
\end{ex}

% nguồn: Ok_TracNghiemLop11_Lan1
\begin{ex}%[0D6K1-3]%[ma:21c7f029]
	Trên  đường tròn lượng giác gốc $A$, cung lượng giác nào có các điểm biểu diễn tạo thành hình vuông?
	\choice 
		 {$\dfrac{k2\pi}{3}$}
		{\True  $\dfrac{k\pi}{2}$}
		{$\dfrac{k\pi}{3}$  }
		{$k\pi $} 
	\loigiai{ 
	\immini{Hình vẽ tham khảo  (hình vẽ bên).
		Hình  vuông $CDEF$ có góc $\widehat{DCE}$ là $45^{\circ}$
		nên góc  ở tâm là $90^{\circ}$ tương ứng $\dfrac{k\pi}{2}.$
		 }{
		\begin{tikzpicture}[>=stealth,scale= 1.3]
				\draw[->] (-1.4,0) -- (1.4,0)  node[above] {$x$};
				 \draw[->] (0,-1.4) -- (0,1.4) node[left] {$y$};
				 \tkzDefPoints{0/0/O,1/0/A,0/1/B}
				\tkzDefPointBy[rotation = center O  angle 30](A) \tkzGetPoint{F}
				 \tkzDefPointBy[symmetry = center O](F)\tkzGetPoint{D}
				 \tkzDefPointBy[rotation = center O angle 90](F) \tkzGetPoint{C}
				\tkzDefPointBy[symmetry =  center O](C)\tkzGetPoint{E}
				\tkzDefPointBy[symmetry =  center O](A)\tkzGetPoint{A'}
				\tkzDefPointBy[symmetry =  center O](B)\tkzGetPoint{B'}
				 \tkzDrawCircle[radius](O,A)
				\tkzDrawSegments (O,D E,C  E,F C,F D,C E,D)
				\tkzLabelPoints[below left](D,B') 
				\tkzLabelPoints[below right](A,E) 
				 \tkzLabelPoints[above right](F,B,O)
				 \tkzLabelPoints[above left](C,A')
			\end{tikzpicture}
		 }
	} 
\end{ex}


%%%%%%%%%%%%%%%% ĐÚNG/SAI %%%%%%%%%%%%%%%%

% nguồn: MaHoaNH_DungSai_Lop10
\begin{ex}%[0D6H1-3]%[ma:abc5e945]
	 Cho $\pi = 3{,} 141 592 653$.
	Các khẳng  định sau đây đúng hay sai?
	\choiceTF 
		{\True Kết quả khi quy  tròn $\pi$ đến hàng phần nghìn là $3{,} 142$}
		{\True Kết quả khi quy tròn $\pi$ đến hàng phần trăm nghìn  là $3{,} 14159$}
		{\True Kết quả khi quy tròn $\pi$ đến hàng phần triệu  là $3{,} 141593$}
		{Kết quả khi quy tròn $\pi$ đến hàng phần trăm triệu là  $3{,} 141592653$}
	\loigiai{ 
		\begin{itemchoice}
			\itemch Kết quả  khi quy tròn $\pi$ đến hàng phần nghìn là $3{,} 142$
			\itemch Kết quả khi quy tròn $\pi$ đến hàng phần trăm nghìn là $3{,} 14159$ 
			\itemch  Kết quả khi quy tròn $\pi$ đến hàng phần triệu là $3{,} 141593$
			\itemch Kết quả khi  quy tròn $\pi$ đến hàng phần trăm triệu là $3{,} 14159265$.
		\end{itemchoice}
	} 
\end{ex}

% nguồn: MaHoaNH_DungSai_Lop10
\begin{ex}%[0D6H1-3]%[ma:905f96c8]
	Cho $\pi = 2{,}718  281 828$.
	Các khẳng  định sau đây đúng hay sai?
	 \choiceTF
		{\True Kết quả khi  quy tròn $\pi$ đến hàng phần trăm là $2{,} 72$}
		{\True Kết quả  khi quy tròn $\pi$ đến hàng phần chục nghìn là $2{,} 7283$}
		{Kết quả  khi quy tròn $\pi$ đến hàng phần mười triệu là $2{,}718 282$}
		{\True Kết quả khi quy tròn $\pi$ đến hàng phần trăm triệu là $2{,}718 281 83$} 
	 \loigiai{
		\begin{itemchoice}
			\itemch Kết  quả khi quy tròn $\pi$ đến hàng phần trăm là $2{,} 72$.
			\itemch  Kết quả khi quy tròn $\pi$ đến hàng phần chục nghìn là $2{,} 7283$.
			 \itemch Kết quả khi quy tròn $\pi$ đến hàng phần mười triệu là $2{,}718 28 18$.
			\itemch Kết quả khi quy tròn $\pi$  đến hàng phần trăm triệu là $2{,}718 281 83$.
		\end{itemchoice}
	 }
\end{ex}


%%%%%%%%%%%%%%%% TỰ LUẬN %%%%%%%%%%%%%%%%

% nguồn: MaHoaNH_TuLuan_Lop10
\begin{bt}%[0D6V1-3]%[ma:dcecc9fe]
	Hình chữ  nhật có các cạnh $x=2 \mathrm{~m} \pm 1 \mathrm{~cm}$, $y=5 \mathrm{~m} \pm 2 \mathrm{~cm}$. Chu vi hình chữ nhật và sai số tương đối của giá trị đó là
	\loigiai{ 
		Chu vi hình chữ nhật  là\\
	\[P_{0}=2\left(x_{0}+y_{0}\right)=2(2+5)=20 m.\]} 
\end{bt}

% nguồn: MaHoaNH_TuLuan_Lop10
\begin{bt}%[0D6V1-3]%[ma:5b0238a3]
	Hãy quy tròn số $\overline{a}=\dfrac{5}{7}=0{,}714285$ đến hàng phần trăm và ước lượng sai  số tương đối.
	 \loigiai{
		Hãy quy tròn số $\overline{a}=\dfrac{5}{7}=0{,}714285$ đến hàng phần trăm và  ước lượng sai số tương đối.\\
		Kết quả quy tròn số  $\overline{a}=\dfrac{5}{7}=0{,}714285$ đến hàng phần trăm là $0{,}71$.\\
		 Sai số tương đối không vượt quá $\dfrac{0{,}005}{071}\approx 0{,}70\%$.
	 }
\end{bt}

% nguồn: MaHoaNH_TuLuan_Lop10
\begin{bt}%[0D6V1-3]%[ma:d42e6da8]
	Độ dài các cạnh của mảnh vườn hình chữ nhật là $x=7{,}8$ m $\pm 2$ cm và $y=25{,}6$  m $\pm 4$ cm. Tìm diện tích (sau khi quy tròn) của mảnh vườn.
	 \loigiai{
		Ta có $\heva{&7{,}78\leq x\leq 7{,}82\\&25{,}56\leq y\leq  25{,}64}\Rightarrow 198{,}8568\leq \overline{S}=xy\leq 200{,}5048$.\\
		Diện tích gần đúng là  $S=7{,}8\cdot 25{,}6=199{,}68$.\\
		 Suy ra $-0{,}8232\leq \overline{S} - S\leq 0{,}8248\Rightarrow\Delta_S=|\overline{S} - S|\leq 0{,}8248$.
		 Suy ra diện tích gần đúng là $S=199{,}68$ m$^2$ với độ chính xác $d=0{,}8248$.
		Nên số quy tròn của $S$ là $200$ m$^2$. 
	 }
\end{bt}
